% ============================================================
% 分片：s06a-superfields  对应原文 PDF 第 75--81 页（印刷页码 75--81）
% 负责章节：6. EXTENSION TO (0,2) MODELS（向 (0,2) 模型的推广）
%           6.1 (0,2) Superfields（(0,2) 超场）
% 译者：t08-s4   审校状态：初稿
% 审校：t06-s33 已完成逐式核对（修正 4 处：0 处公式、4 条批注）
% ============================================================
% 记法说明（照原文，不改符号）：
%   • 规范协变导数为花体 $\mathcal{D}$；超空间微分算符为普通斜体 $D$。原文 (6.1)
%     的算符记 $Q$、(6.2) 记 $D$，两者只差 $\partial_0+\partial_1$ 项的符号。
%   • 原文以粗体 $\mathbb{C}$ 等排版处按 GLOSSARY/SPEC 4.5 统一为 $\mathbb{C}$；
%     微分与积分测度按原文直立排印，写作 $\mathrm{d}$；"h.c." 按原文直立。
%   • 原文脚注以 $\star$ 标注，本译本改用 \footnote，编号由 LaTeX 重新生成。
%   • 第 6.1 节共 34 个带号公式：(6.1)--(6.34)，全部锁号，无一遗漏。
%   • 分片止于第 81 页末（式 (6.34) 所在段落结束）；6.2 的标题及正文归 s06b 分片。

% ==== 原文 PDF 第 75 页（印刷页 75）====
\section{向 $(0,2)$ 模型的推广}

\begin{keytakeaway}
\textbf{核心物理图像：}
$(0,2)$ 超对称即两个右行、零个左行超对称。把第 3 节的线性 $\sigma$ 模型构造逐项移植到 $(0,2)$ 超空间，便得到 $(0,2)$ 版本的 Calabi--Yau/Landau--Ginzburg 对应：原始模型的场按 $(2,2)\to(0,2)$ 分裂，超势 $W$ 被一对由全纯函数 $E_a$、$J^a$ 描述的数据取代，其玻色势由 $D$ 项、$|E|^2$、$|J|^2$ 三部分组成。
\par\textbf{数学工具：}
$(0,2)$ 超空间与手征/矢量超场约束、Wess--Zumino 规范、超空间积分 $\int\mathrm{d}^2y\,\mathrm{d}\theta^+\,\mathrm{d}\bar\theta^+$ 与 $\int\mathrm{d}^2y\,\mathrm{d}\theta^+|_{\bar\theta^+=0}$、$(2,2)$ 多重态在 $\theta^-=\bar\theta^-=0$ 处的截断、全纯函数与规范不变性。
\par\textbf{主要结论：}
规范多重态由 $v_\alpha,\lambda_-,\bar\lambda_-,D$ 组成，手征多重态为 $\Phi$，Fermi 多重态 $\Lambda_-$ 满足 $\bar{\mathcal{D}}_+\Lambda_-=\sqrt2E$；作用量 $L=L_{gauge}+L_{ch}+L_F+L_{D,\theta}+L_J$ 消去辅助场后给出势 (6.27)。$(2,2)$ 模型约化时 $J^i=\partial W/\partial\Phi'_i$、$E_i=iQ_i\sqrt2\Sigma'\Phi'_i$，而 $E_iJ^i=0$ 恰是 $W$ 规范不变性的推论，这为第 6.2 节的 $(0,2)$ 模型族提供了出发点。
\end{keytakeaway}

具有 $(0,2)$ 超对称（即两个右行、零个左行超对称，与上文所研究的 $(2,2)$ 模型相对）的模型在唯象上很重要。在弦论紧化到四维的背景下，它们自然地给出 $SU(5)$ 或 $SO(10)$、而非 $E_6$，作为有效的大统一规范群。这些模型素有“比 $(2,2)$ 模型难研究得多”的名声，因为 Gepner 可解模型、Landau--Ginzburg 对应等的类比物尚不为人所知。本节将构造 Landau--Ginzburg 对应的 $(0,2)$ 版本。这不过是一项直截了当的工作：弄清相应的 $(0,2)$ 超场的结构，然后重复第 3 节的手续。

\subsection{$(0,2)$ 超场}

我们将在此 $(0,2)$ 超空间中工作：玻色坐标为 $y^\alpha$，$\alpha=1,2$，费米坐标为 $\theta^+$、$\bar\theta^+$。\footnote{此前已有与下文相关的讨论 [49,50,51]，但下文有些论点是新的，另一些则是为完整起见而重述的。} 超对称生成元为

% [原文疑误: 原文说玻色坐标 $y^\alpha$ 的指标 $\alpha=1,2$，但 (6.1)、(6.2) 及其后通篇用的是 $\partial/\partial y^0$、$\partial/\partial y^1$，且 (6.5)、(6.6) 用 $\mathcal{D}_0\pm\mathcal{D}_1$。故 $\alpha=1,2$ 应是 $\alpha=0,1$ 的笔误（二维时空指标从 0 起），译文照原文保留 $y^\alpha,\alpha=1,2$，建议审校时按 $\alpha=0,1$ 理解。]
\begin{equation}
\begin{aligned} Q_+ &= \frac{\partial}{\partial\theta^+}+i\bar\theta^+\left(\frac{\partial}{\partial y^0}+\frac{\partial}{\partial y^1}\right)\\ \bar Q_+ &= -\frac{\partial}{\partial\bar\theta^+}-i\theta^+\left(\frac{\partial}{\partial y^0}+\frac{\partial}{\partial y^1}\right). \end{aligned}\tag{6.1}
\end{equation}
这些算符与以下算符对易：
\begin{equation}
\begin{aligned} D_+ &= \frac{\partial}{\partial\theta^+}-i\bar\theta^+\left(\frac{\partial}{\partial y^0}+\frac{\partial}{\partial y^1}\right)\\ \bar D_+ &= -\frac{\partial}{\partial\bar\theta^+}+i\theta^+\left(\frac{\partial}{\partial y^0}+\frac{\partial}{\partial y^1}\right), \end{aligned}\tag{6.2}
\end{equation}
% ==== 原文 PDF 第 76 页（印刷页 76）====
\label{op:76}
% [跨页接续: 原文第 75 页/第 76 页；本句跨页，按原文不切分句子]
而后者用于构造拉格朗日量。

% [Tutor Note: 由 (6.1)、(6.2) 直接算得 $\{Q_+,\bar Q_+\}=-2i(\partial_0+\partial_1)$、$\{D_+,\bar D_+\}=+2i(\partial_0+\partial_1)$，且同手征的 $\{Q_+,D_+\}=\{\bar Q_+,\bar D_+\}=0$（混合手征的反对易子不为零，正比于 $\partial_0+\partial_1$）——同手征的 $Q$ 与 $D$ 恰好只差 $\partial_0+\partial_1$ 项的符号，故互相反对易。$\partial_0+\partial_1$ 是沿 $y^0+y^1$ 光锥方向的导数，这解释了为何 $(0,2)$ 模型只有“右行”超对称：仅 $\theta^+$（手征指标 $+$）参与配对。易错点：本节的 $\bar\theta^+$ 与 $\theta^+$ 互为共轭，$\bar\theta^+$ 上的横线不是狄拉克共轭指标，切勿与 (2,2) 中的 $\bar\theta^{\dot\alpha}$ 混同。]

作者在此并不声称描述了所有可能的 $(0,2)$ 模型，而只描述那些能用某些类型的 $(0,2)$ 超场方便地描述的模型。

\subsubsection*{规范多重态}

我们想在超空间中引入规范场。规范协变导数将记为 $\mathcal{D}_+$、$\bar{\mathcal{D}}_+$ 以及 $\mathcal{D}_\alpha=D/Dy^\alpha$。我们假定超空间规范场满足约束
\begin{equation}
\begin{aligned} \mathcal{D}_+^2 &= \bar{\mathcal{D}}_+^2=0\\ \mathcal{D}_+\bar{\mathcal{D}}_++\bar{\mathcal{D}}_+\mathcal{D}_+ &= 2i(\mathcal{D}_0+\mathcal{D}_1). \end{aligned}\tag{6.3}
\end{equation}
（例如，这些方程允许稍后引入的 $(0,2)$ 手征超场存在。）前两个方程意味着
\begin{equation}
\begin{aligned} \mathcal{D}_+ &= e^{-\Psi}D_+e^{\Psi}\\ \bar{\mathcal{D}}_+ &= e^{\bar\Psi}\bar D_+e^{-\bar\Psi}, \end{aligned}\tag{6.4}
\end{equation}
其中 $\Psi$ 是李代数值的函数。通过规范变换可以假定 $\Psi$ 为实。也可以把 $\Psi$ 中与 $\theta^+$ 无关、或与 $\bar\theta^+$ 无关的项规范掉。特别地，可以取 Wess--Zumino 规范的类比，其中 $\Psi=\theta^+\bar\theta^+(v_0+v_1)$，而 $v_0+v_1$ 只是 $y^\alpha$ 的函数。在此部分规范固定下，
\begin{equation}
\begin{aligned} \mathcal{D}_0+\mathcal{D}_1 &= \partial_0+\partial_1+i(v_0+v_1)\\ \mathcal{D}_+ &= \frac{\partial}{\partial\theta^+}-i\bar\theta^+(\mathcal{D}_0+\mathcal{D}_1)\\ \bar{\mathcal{D}}_+ &= -\frac{\partial}{\partial\bar\theta^+}+i\theta^+(\mathcal{D}_0+\mathcal{D}_1). \end{aligned}\tag{6.5}
\end{equation}
超空间规范场的其余部分是
\begin{equation}
\mathcal{D}_0-\mathcal{D}_1=\partial_0-\partial_1+iV,\tag{6.6}
\end{equation}
% ==== 原文 PDF 第 77 页（印刷页 77）====
\label{op:77}
% [跨页接续: 原文第 76 页/第 77 页；本句跨页，按原文不切分句子]
其中某个函数 $V$ 可展开为
\begin{equation} V=v_0-v_1-2i\theta^+\bar\lambda_--2i\bar\theta^+\lambda_-+2\theta^+\bar\theta^+D.\tag{6.7}
\end{equation}

% [Tutor Note: (6.3) 在平直超空间中前两式自动成立（$\partial_\theta^2=\bar\theta^2=0$、$\{\partial_\theta,\bar\theta\}=0$），成为规范协变形式后它们才是对超空间规范场的真实约束（即超空间挠率为零）。第三式则定义了 $\mathcal{D}_0+\mathcal{D}_1$ 中出现的组合 $v_0+v_1$；注意 $\mathcal{D}_0+\mathcal{D}_1=\partial_0+\partial_1+i(v_0+v_1)$ 正是“沿 $y^0+y^1$ 方向的协变导数”，与 (6.1)--(6.2) 中只出现 $\partial_0+\partial_1$ 一致。]

% [Tutor Note: 由 (6.6) 的 $\mathcal{D}_0-\mathcal{D}_1=\partial_0-\partial_1+iV$ 与 (6.7) 可读出分量计数：$\Psi$ 的 $\theta^+\bar\theta^+$ 分量给出两个玻色分量 $v_0,v_1$（即光锥分量 $v_0\pm v_1$），$V$ 的 $\theta^+$、$\bar\theta^+$、$\theta^+\bar\theta^+$ 分量给出两个费米子 $\lambda_-,\bar\lambda_-$ 与一个辅助场 $D$；系数 $2i$、$2$ 由 (6.5) 的规范协变形式定出。]

$v_\alpha$、$\lambda_-$、$\bar\lambda_-$ 与 $D$ 的变换律可由这些公式推出（利用超对称生成元 (6.2)，并辅以保持 (6.5) 形式的规范变换），它们恰好是 $(2,2)$ 规范多重态变换律 (2.12) 的 $(0,2)$ 截断——只是某些场缺失了。缺失的场构成一个另立的 $(0,2)$ 多重态，稍后将予以辨认。

% [存疑: 原文此处说“利用超对称生成元 (6.2)”，但 (6.2) 是上面刚定义为与超荷对易的微分算符 $D$、$\bar D$，而超对称生成元是 (6.1) 的 $Q$、$\bar Q$。疑为 (6.1) 之误；按 SPEC 照录原文并在此标注，供审校核定。]

基本的规范不变场强是 $\Upsilon=[\bar{\mathcal{D}}_+,\mathcal{D}_0-\mathcal{D}_1]$。$(0,2)$ 规范多重态的自然超空间作用量为
\begin{equation} L_{gauge}=\frac{1}{8e^2}\int\mathrm{d}^2y\,\mathrm{d}\theta^+\,\mathrm{d}\bar\theta^+\,\operatorname{Tr}\bar\Upsilon\Upsilon.\tag{6.8}
\end{equation}
将其在分量上算出（为简单起见取 $U(1)$ 情形）给出
\begin{equation} \frac{1}{e^2}\int\mathrm{d}^2y\left(\frac{1}{2}v_{01}^2+i\bar\lambda_-(\partial_0+\partial_1)\lambda_-+\frac{1}{2}D^2\right).\tag{6.9}
\end{equation}

% [Tutor Note: (6.9) 的归一化可用于核对 (6.8) 中的 $1/8e^2$：在 $U(1)$ 情形 $\operatorname{Tr}$ 为恒等，$\Upsilon$ 的 $\theta^+$ 分量给出 $2i(D-iv_{01})$ 一类的组合，$\int\mathrm{d}\theta^+\mathrm{d}\bar\theta^+$ 取 $\theta^+\bar\theta^+$ 分量后正好给出 $v_{01}^2/2+D^2/2$，故前置系数为 $1/8e^2$。易错点：$v_{01}=\partial_0v_1-\partial_1v_0$ 是二维场强，与 (6.7) 中 $v_0-v_1$ 的符号约定直接相关。]

\subsubsection*{手征多重态}

需要考虑两类物质多重态。其一是某个规范群表示中的玻色场 $\Phi$，满足
\begin{equation}
\bar{\mathcal{D}}_+\Phi=0.\tag{6.10}
\end{equation}
我们称之为 $(0,2)$ 手征多重态。手征多重态具有 $\theta$ 展开
\begin{equation}
\Phi=\phi+\sqrt2\theta^+\psi_+-i\theta^+\bar\theta^+(D_0+D_1)\phi.\tag{6.11}
\end{equation}
（此处 $D_\alpha$ 当然是 $\theta^+=\bar\theta^+=0$ 处的“普通”规范协变导数。）若 $\Phi$ 是电荷为 $Q$ 的场，并与上述阿贝尔规范多重态相互作用，
% ==== 原文 PDF 第 78 页（印刷页 78）====
\label{op:78}
% [跨页接续: 原文第 77 页/第 78 页；本句跨页，按原文不切分句子]
则在分量上，作用量为
\begin{equation}
\begin{aligned} L_{ch} &= -\frac{i}{2}\int\mathrm{d}^2y\,\mathrm{d}^2\theta\;\bar\Phi(\mathcal{D}_0-\mathcal{D}_1)\Phi=\int\mathrm{d}^2y\Big(-|D_\alpha\phi|^2+\bar\psi_+i(D_0-D_1)\psi_+\\ &\qquad\qquad-iQ\sqrt2\bar\phi\lambda_-\psi_++iQ\sqrt2\bar\psi_+\bar\lambda_-\phi+QD\bar\phi\phi\Big). \end{aligned}\tag{6.12}
\end{equation}

% [Tutor Note: (6.12) 中两个汤川项互为厄米共轭（相差费米子交换的符号），故整个 $L_{ch}$ 为实；$QD\bar\phi\phi$ 在消去 $D$ 后给出 $D$ 项势（见 (6.27) 第一项）。易错点：$|D_\alpha\phi|^2$ 中的 $D_\alpha$ 是 $\theta=0$ 处的规范协变导数，与 (6.11) 里同名的 $D_0,D_1$ 一致，而非超空间微分算符 $D_+$。]

\subsubsection*{Fermi 多重态}

另一类物质多重态是某个规范群表示中的反对易、负手征旋量场 $\Lambda_-$，满足
\begin{equation}
\bar{\mathcal{D}}_+\Lambda_-=\sqrt2E,\tag{6.13}
\end{equation}
其中 $E$ 是某个满足下式的超场：
\begin{equation}
\bar{\mathcal{D}}_+E=0.\tag{6.14}
\end{equation}
我们将称 $\Lambda_-$ 为一个 Fermi 多重态（fermi multiplet）。Fermi 多重态的 $\theta$ 展开为
\begin{equation}
\Lambda_-=\lambda_--\sqrt2\theta^+G-i\theta^+\bar\theta^+(D_0+D_1)\lambda_--\sqrt2\bar\theta^+E.\tag{6.15}
\end{equation}

% [Tutor Note: (6.15) 中 $-\sqrt2\bar\theta^+E$ 一项由约束 (6.13) 唯一确定：作用 $\bar{\mathcal{D}}_+$ 时它给出 $\sqrt2E$，其余各项给出的 $\bar\theta$ 无关项或 $\bar\theta^+$ 项都因 $\bar{\mathcal{D}}_+$ 的显式形式（见 (6.5)；平直情形见 (6.2)）而被抵消。$G$ 是 Fermi 多重态的辅助场，在 (6.18) 中仅以 $|G|^2$ 出现，消去后不产生势。]

$E$ 自身也有 theta 展开。在 $E=E(\Phi_i)$ 是某些手征超场 $\Phi_i$ 的全纯函数这一重要情形（这些超场的展开如 (6.11) 所示），有
\begin{equation} E(\Phi_i)=E(\phi_i)+\sqrt2\theta^+\frac{\partial E}{\partial\phi_i}\psi_{+,i}-i\theta^+\bar\theta^+(D_0+D_1)E(\phi_i).\tag{6.16}
\end{equation}
Fermi 多重态的自然作用量为
\begin{equation} L_F=-\frac{1}{2}\int\mathrm{d}^2y\,\mathrm{d}^2\theta\,\bar\Lambda_-\Lambda_-.\tag{6.17}
\end{equation}
若 $E$ 如 (6.16) 所示，则分量展开为
\begin{equation}
\begin{aligned} L_F=\int\mathrm{d}^2y\bigg(&i\bar\lambda_-(D_0+D_1)\lambda_-+|G|^2-|E(\phi_i)|^2\\ &-\left(\bar\lambda_-\frac{\partial E}{\partial\phi_i}\psi_{+,i}+\frac{\partial\bar E}{\partial\bar\phi_i}\bar\psi_{+,i}\lambda_-\right)\bigg). \end{aligned}\tag{6.18}
\end{equation}

% [Tutor Note: (6.18) 中 $-|E|^2$ 即 $(0,2)$ 版的“$F$ 项势”，而 $|G|^2$ 只是辅助场项；对 $G$ 变分得 $G=0$，故 $G$ 不产生势。注意与 (6.17) 的符号约定相应配：$\int\mathrm{d}^2\theta\,\bar\Lambda_-\Lambda_-$ 取 $\theta^+\bar\theta^+$ 分量时产生 $i\bar\lambda(\partial_0+\partial_1)\lambda$ 与 $-|E|^2$，二者符号相反是 (0,2) 超空间的标准结果。]

% ==== 原文 PDF 第 79 页（印刷页 79）====
\label{op:79}
\subsubsection*{作用量中的其他项}

作用量中的其他项将具有形式
\begin{equation}
\int\mathrm{d}^2y\,\mathrm{d}\theta^+\,(\ldots)\big|_{\bar\theta^+=0}+\mathrm{h.c.},\tag{6.19}
\end{equation}
其中 $\ldots$ 是某个被 $\bar{\mathcal{D}}_+$ 湮灭的反对易超场。这些项无法写成对整个超空间的积分，却格外重要。

例如，规范场强 $\Upsilon$ 满足 $\bar{\mathcal{D}}_+\Upsilon=0$。于是在阿贝尔情形下可以写出
\begin{equation}
\begin{aligned} L_{D,\theta} &= \frac{t}{4}\int\mathrm{d}^2y\,\mathrm{d}\theta^+\,\Upsilon\big|_{\bar\theta^+=0}+\mathrm{h.c.}\\ &= \frac{it}{2}\int\mathrm{d}^2y\,(D-iv_{01})-\frac{i\bar t}{2}\int\mathrm{d}^2y\,(D+iv_{01}). \end{aligned}\tag{6.20}
\end{equation}

% [Tutor Note: (6.20) 中 $\int\mathrm{d}\theta^+\Upsilon|_{\bar\theta^+=0}$ 取出的是 $D$ 与 $v_{01}$ 的组合：按 (6.20) 第二个等号反推，$\int\mathrm{d}\theta^+\Upsilon|_{\bar\theta^+=0}=2i(D-iv_{01})$（归一化取决于 (6.8) 的约定）。$t$ 是复化 FI 参数（第 3 节的 $t$，与 $r$ 及二维 $\theta$ 角相关，见第 75 页开头与第 73 页的讨论）；$t$ 与 $\bar t$ 分别乘以 $D-iv_{01}$ 与 $D+iv_{01}$，正是二维规范理论中 FI 项与 $\theta$ 项的标准组合。]

作为另一个主要例子，令 $\Lambda_{-,a}$ 为某些 Fermi 超场，满足
\begin{equation}
\bar{\mathcal{D}}_+\Lambda_{-,a}=\sqrt2E_a,\tag{6.21}
\end{equation}
$E_a$ 为某些手征超场。再令 $J^a$ 为某些手征超场，满足
\begin{equation}
E_aJ^a=0.\tag{6.22}
\end{equation}
那么
\begin{equation}
\bar{\mathcal{D}}_+(\Lambda_{-,a}J^a)=0.\tag{6.23}
\end{equation}
于是我们可以在作用量中引入另一项；它是超势的 $(0,2)$ 类比：
\begin{equation} L_J=-\frac{1}{\sqrt2}\int\mathrm{d}^2y\,\mathrm{d}\theta^+\,\Lambda_{-,a}J^a\big|_{\bar\theta^+=0}-\mathrm{h.c.}\tag{6.24}
\end{equation}
若假定 $E_a$ 与 $J^a$ 是手征超场 $\Phi_i$ 的全纯函数，
% ==== 原文 PDF 第 80 页（印刷页 80）====
\label{op:80}
% [跨页接续: 原文第 79 页/第 80 页；本句跨页，按原文不切分句子]
便得到
\begin{equation} L_J=-\int\mathrm{d}^2y\left(G_aJ^a(\phi_i)+\lambda_{-,a}\psi_{+,i}\frac{\partial J^a}{\partial\phi_i}\right)-\mathrm{h.c.}\tag{6.25}
\end{equation}

% [Tutor Note: (6.22) 的必要性可直接验证：$J^a$ 手征（$\bar{\mathcal{D}}_+J^a=0$）、$\Lambda_{-,a}$ 满足 (6.21)，故 $\bar{\mathcal{D}}_+(\Lambda_{-,a}J^a)=(\bar{\mathcal{D}}_+\Lambda_{-,a})J^a=\sqrt2E_aJ^a$；令其为零即得 (6.22)，从而 (6.23) 成立，$\Lambda_{-,a}J^a$ 才可作 $\theta^+$ 积分。这就是 $(0,2)$ 超势项存在的充要条件。]

综合所有这些，取拉格朗日量为
\begin{equation}
L=L_{gauge}+L_{ch}+L_F+L_{D,\theta}+L_J.\tag{6.26}
\end{equation}
消去辅助场 $D$ 与 $G_a$ 后，玻色势为
\begin{equation} U(\phi_i)=\frac{e^2}{2}\left(\sum_iQ_i|\phi_i|^2-r\right)^2+\sum_a|E_a|^2+\sum_a|J^a|^2.\tag{6.27}
\end{equation}

% [Tutor Note: $D$ 的方程给出 $D=-e^2(\sum_iQ_i|\phi_i|^2-r)$，代回 (6.20) 与 (6.12) 的 $QD\bar\phi\phi$ 即得 (6.27) 第一项；$G_a$ 的方程给出 $G_a=\bar J^a$（即 $\bar G_a=J^a$），代回 (6.25) 与 (6.18) 的 $|G|^2$ 即得 $\sum_a|J^a|^2$。故 (6.27) 恰是 $(2,2)$ 势 $\frac{e^2}{2}(\sum Q_i|\phi_i|^2-r)^2+\sum_i|\partial W/\partial\phi_i|^2$ 的 $(0,2)$ 类比：$|E|^2$ 与 $|J|^2$ 分担了 $|\partial W|^2$ 的角色。]

$(0,2)$ 版的 C.Y./L.G. 对应将按如今已熟悉的方式、通过研究真空结构随 $r$ 的依赖关系而得到。

\subsubsection*{$(2,2)$ 多重态的约化}

现在讨论 $(2,2)$ 多重态到 $(0,2)$ 多重态的约化。

首先考虑 $(2,2)$ 规范多重态，它在 Wess--Zumino 规范下由标量函数 $V(y,\theta^\pm,\bar\theta^\pm)$ 描述，其展开由式 (2.11) 给出。该多重态的一部分组成了上面描述的 $(0,2)$ 规范多重态。剩余的场是 $\sigma$、$\lambda_+$ 与 $\bar\lambda_+$，它们组成一个 $(0,2)$ 手征多重态。事实上，这个多重态正是 $\Sigma'=\Sigma|_{\theta^-=\bar\theta^-=0}$，其中 $\Sigma=\sigma+\ldots$ 是 $(2,2)$ 规范多重态在式 (2.15) 中引入的规范不变场强。\footnote{为把 $\Sigma'$ 的 $\theta$ 展开写成 $(0,2)$ 手征多重态的标准形式 (6.11)，必须在 $\bar\lambda_+$ 中吸收一个因子 $i$。}

再考虑 $(2,2)$ 手征多重态 $\Phi$。$\Phi$ 在 $(0,2)$ 超对称下分解为两个多重态。首先是一个 $(0,2)$ 手征多重态
\begin{equation}
\Phi'=\Phi\big|_{\theta^-=\bar\theta^-=0}.\tag{6.28}
\end{equation}
% ==== 原文 PDF 第 81 页（印刷页 81）====
\label{op:81}
其次，令
\begin{equation}
\Lambda_-=\frac{1}{\sqrt2}\mathcal{D}_-\Phi\big|_{\theta^-=\bar\theta^-=0}.\tag{6.29}
\end{equation}
则
\begin{equation}
\bar{\mathcal{D}}_+\Lambda_-=\frac{1}{\sqrt2}\{\bar{\mathcal{D}}_+,\mathcal{D}_-\}\Phi\big|_{\theta^-=\bar\theta^-=0}.\tag{6.30}
\end{equation}
例如，若规范群为 $U(1)$、$\Phi$ 的电荷为 $Q$，这相当于
\begin{equation}
\bar{\mathcal{D}}_+\Lambda_-=2iQ\Sigma'\Phi'.\tag{6.31}
\end{equation}
所以 $\Lambda_-$ 是一个 Fermi 多重态，其
\begin{equation}
E=iQ\sqrt2\Sigma'\Phi'.\tag{6.32}
\end{equation}

% [Tutor Note: 由 (6.30) 到 (6.31) 的换算要用到 $(2,2)$ 规范多重态的场强 $\Sigma$ 的定义（式 (2.15)），并注意 $\{\bar{\mathcal{D}}_+,\mathcal{D}_-\}\Phi=2\sqrt2iQ\Sigma'\Phi'$ 这一代数关系；$1/\sqrt2$ 与 $\sqrt2$ 抵消后得 $2iQ$。若把 (6.15) 的标准形式与 (6.29) 对照，还需按第 80 页脚注在 $\bar\lambda_+$ 中吸收因子 $i$；(6.31) 的 $2iQ\Sigma'\Phi'$ 再经 (6.13) 的 $\sqrt2$ 因子，即给出 (6.32) 的 $iQ\sqrt2\Sigma'\Phi'$。]

最后，我们想要 $(2,2)$ 模型超势的 $(0,2)$ 约化。考虑一个含手征超场 $\Phi_i$ 与超势 $W(\Phi_i)$ 的 $(2,2)$ 模型。在 $(0,2)$ 超对称下，如刚才所见，$\Phi_i$ 分裂为 $(0,2)$ 手征超场 $\Phi'_i$ 与 Fermi 超场 $\Lambda_{-,i}$。后者满足 $\bar{\mathcal{D}}_+\Lambda_{-,i}=\sqrt2E_i$，其中由 (6.32)，
\begin{equation}
E_i=iQ_i\sqrt2\Sigma'\Phi'_i.\tag{6.33}
\end{equation}
在 $(2,2)$ 超对称中，手征超场的标量耦合与汤川耦合由单个函数 $W$ 决定；而在 $(0,2)$ 超对称中，我们必须指定一组函数 $J^i$，每个 $\Lambda_{-,i}$ 对应一个，并要求 $J^iE_i=0$，其中 $E_i$ 如上。由 $(2,2)$ 模型约化得到的 $J^i$ 正是
\begin{equation}
J^i=\frac{\partial W}{\partial\Phi'_i}.\tag{6.34}
\end{equation}
这可以通过把 $(2,2)$ 模型中由超势导出的耦合的式 (2.21) 与由 $J$ 确定的耦合的式 (6.25) 相比较而得到。取 (6.34) 的 $J^i$ 与 (6.33) 的 $E_i$，方程 $E_iJ^i=0$ 是 $W$ 的规范不变性的推论。

% [Tutor Note: 验证 $E_iJ^i=0$：$E_iJ^i=i\sqrt2\Sigma'\sum_iQ_i\Phi'_i\partial W/\partial\Phi'_i$。$W$ 规范不变（电荷为零）意味着 $\sum_iQ_i\Phi'_i\partial W/\partial\Phi'_i=0$（对每个单项式，其总电荷为零，故欧拉型恒等式成立），于是 $E_iJ^i=0$ 正是 $W$ 规范不变性的直接结果。这也说明：$(0,2)$ 模型把 $(2,2)$ 中“$W$ 的规范不变性”这一约束换成了“$E_iJ^i=0$”这一更弱的代数条件，从而给出更大的模型族（第 6.2 节即以此为出发点）。]
